% !TeX root = ../../Book4.tex
% 纲要标题：### 17.4 映射锥的函子性与选择
% 纲要标签：b4:sec:1604
% 纲要位置：计划纲要.md，第 2913 行。

\section{映射锥的函子性与选择}
\label{b4:sec:1604}

映射锥在复形及其严格交换方块组成的箭头范畴上构成函子。只保留同伦类及普通三角中的交换方块时，补全锥映射所用的同伦没有被指定，因而这些数据本身不决定典范的锥映射。


% 纲要内容（保留纲要核验项目）：
%
% * 三角范畴中锥的选择
% * 同构不等于典范选择
% * DG 增强的动机
%
% ---
%

\subsection{三角范畴中锥的选择}
\label{b4:subsec:1604:01}

在复形层面，选定 \(f\) 后，\(\operatorname{Cone}(f)=B\oplus A[1]\) 有明确公式。若方块 \(bf=f'a\) 严格交换，则 \((y,x)\mapsto(by,ax)\) 给出两个锥之间的映射，而且保持复合。因此严格交换方块构成的箭头范畴上确实有锥函子。

\ProseParagraphBreak
转到同伦范畴，方块只要求 \(bf-f'a=\mathrm{d}h+h\mathrm{d}\)。构造锥之间的映射便需要选 \(h\)，得到 \((y,x)\mapsto(by+hx,ax)\)。换一个同伦，结果未必同伦。这正是 (TR3) 只要求存在补全的原因：先忘掉同伦本身，再仅凭四条边，通常无法恢复原先的锥映射。

\subsection{同构与典范选择}
\label{b4:subsec:1604:02}

取域 \(k\)，令 \(X=k[-1]\)、\(Y=k[0]\)，考虑零映射 \(X\to Y\)。其锥是次数零上的 \(k\oplus k\)，三角为
\[
 X\xrightarrow{0}Y\longrightarrow Y\oplus X[1]\longrightarrow X[1].
\]
对每个 \(\lambda\in k\)，矩阵
\[
 c_\lambda=\begin{pmatrix}1&\lambda\\0&1\end{pmatrix}
\]
都与包含和投影相容，因此 \((1_X,1_Y,c_\lambda)\) 是三角自同构。不同 \(\lambda\) 在同伦范畴中仍不同，因为次数零复形之间没有能改变映射的非零链同伦。

\ProseParagraphBreak
这个例子中，固定前两个分量仍得到不同的第三分量，因而相容同构没有唯一性。若要让锥映射随方块函子性地变化，还须指定相容的选择，并检验这些选择保持恒等态射与复合。

\subsection{DG 增强的动机}
\label{b4:subsec:1604:03}

要保留上面的 \(h\)，可以不把 Hom 立即压缩成同伦类，而保存整个 Hom 复形。零次余圈是复形映射，次数负一的元素是候选同伦，等式 \(\mathrm{d}h=f-g\) 记录两映射为何同伦。进一步的负次数还记录同伦之间的相容关系。

\ProseParagraphBreak
这给出\term[DG-zengqiang]{DG 增强}[DG enhancement]的动机：一个普通范畴之外，还保留态射复形及与微分相容的复合。取这些复形的零次上同调后，才回到同伦范畴。\chapref{b4:ch:20}将给出正式定义和实例。
